\section{\systemname Runtime}
\label{sec:method}

\subsection{System Overview}

\systemname is a general-purpose agentic runtime built around four model-driven roles
and three system planes. The \emph{control plane} anchors the campaign and schedules
work; the \emph{execution plane} performs one bounded mission against real tools
and artifacts; and the \emph{record plane} stores what happened without deciding
whether the work is scientifically complete. The distinction is operational:
control owns scheduling, execution owns work and mission review, and the record plane owns
the immutable record but no completion decision.

The four roles divide research authority as summarized in
Table~\ref{tab:roles}. The Manager spans the campaign rather than one execution
round. The Planner authors future work and the Engineer performs it. Mission-level
completion has an explicit source: allowed low-risk bounded work may use Engineer
self-review, while vertical policy, stage-closing tasks, or the Engineer's own
request require an independent Reviewer. The interfaces are structured objects
rather than untyped prose, which makes ownership and recovery behavior explicit.

\begin{table}[H]
\centering
\small
\begin{tabularx}{\linewidth}{@{}p{1.6cm} p{4.0cm} X@{}}
\toprule
\rowcolor{papersand}
\textbf{Role} & \textbf{Primary responsibility} & \textbf{Authoritative output} \\
\midrule
Manager & Commit the objective, lifetime, route, and campaign stage & Campaign division and stage transition \\
Planner & Decompose the current research state into bounded tasks and dependencies & Task specifications, plan status, and stage checklist updates \\
Engineer & Modify artifacts, invoke tools, and run experiments for one round & Execution result, artifacts, measurements, handoff, and structured \code{review=skip|required} selection \\
Reviewer & Independently inspect artifacts and execution records when required or requested & \code{done}, \code{continue}, or \code{blocked}, plus grounded next action and retained-state edits \\
\bottomrule
\end{tabularx}
\caption{Role separation in \systemname. The Engineer executes and may self-review
allowed low-risk bounded work; mandatory or requested independent review belongs to
the Reviewer.}
\label{tab:roles}
\end{table}

\begin{paperbox}{Algorithm 1: Campaign scheduling and the reviewed mission loop}
\small
\begin{tabbing}
\quad\=\quad\=\quad\=\quad\=\kill
\textbf{Input:} standing intent $\iota$, contract $K_t$, runtime state $H_t$, backlog $B$, budget $b$ \\
$c \leftarrow \operatorname{ManagerCommit}(\iota,K_t)$; persist campaign identity \\
\textbf{while} $b$ remains and $c$ is not complete \textbf{do} \\
\> \textit{// campaign level: the mission boundary is fixed here} \\
\> $N \leftarrow \operatorname{StableTopological}(\operatorname{Planner}(B,H_t,c,K_t))$ \\
\> \textbf{for} each bounded mission $q_0 \in N$, enqueued one at a time \textbf{do} \\
\>\> $L \leftarrow \operatorname{ExposeSkillLibraries}(q_0)$; \textit{paths only, retrieval is the agent's} \\
\>\> $\Gamma \leftarrow [\,]$; $q \leftarrow q_0$ \\
\>\> \textbf{repeat} \textit{// mission level: $q_0$ is now fixed for every round below} \\
\>\>\> $(y,e,d) \leftarrow \operatorname{Engineer}(q,H_t,L)$ \\
\>\>\> \textbf{if} $d=\code{skip}$ and $\operatorname{SelfReviewAllowed}(q)$ \\
\>\>\>\> $r \leftarrow \operatorname{EngineerSelfReview}(q,y,e)$ \\
\>\>\> \textbf{else} $r \leftarrow \operatorname{Reviewer}(q,y,e)$ \textit{; scoped to this round} \\
\>\>\> append $(q,y,e,r,r.\code{source})$ to $\Gamma$ \\
\>\>\> \textbf{if} $r=\code{continue}$ \textbf{then} $q \leftarrow \operatorname{AdaptAfterRejections}(\Gamma)$ \\
\>\> \textbf{until} $r\in\{\code{done},\code{blocked},\code{paused}\}$ or a governance threshold fires \\
\>\> $(\code{status},\code{reason}) \leftarrow \operatorname{PreSettlementGuard}(\Gamma,\code{status},\code{reason})$ \\
\>\> $\tau \leftarrow \operatorname{TraceProjection}(\Gamma)$; \quad $E \leftarrow \operatorname{AdmitResult}(\tau)$ \\
\>\> $H_{t+1} \leftarrow \operatorname{SettleMissionOutcome}(H_t,\tau,E)$ \textit{; learning happens here} \\
\> \textbf{if} evidence proposes a contract change $K'_t$ \textbf{then} \\
\>\> $(K_{t+1},\rho_t) \leftarrow \operatorname{ReviseContract}(K_t,K'_t,\code{by},\code{confirmation})$ \\
\> $t \leftarrow t+1$; refill $B$ when required \\
\textbf{end while} \\
\textbf{Output:} accepted artifacts, inspectable trajectory, admitted contract, and updated runtime state
\end{tabbing}
\end{paperbox}

The algorithm above separates two levels that are easy to conflate. The
Planner runs at the campaign level and fixes the mission boundary $q_0$; the round
loop beneath it varies $q$ only through \code{AdaptAfterRejections}, which reacts to
Reviewer rejections \emph{within} that boundary. No component inside the round loop
can rewrite $q_0$. Section~\ref{sec:endogenous-harness} treats the consequence of
that asymmetry as a first-class failure mode rather than an implementation detail.

Three further details matter because they differ from the idealized loop.
\emph{Skill libraries are exposed, not injected}: the runtime publishes library paths
and the acting agent performs its own retrieval, so no matcher decides in advance
which prior experience is relevant. \emph{Termination is governed by named
thresholds} rather than by Reviewer verdicts alone: a maximum round count, a
no-progress threshold, a soft round limit, a hard escalation count, and a backend
failure threshold each end or escalate a mission, and a pre-settlement guard may
override the recorded status before anything is learned. \emph{Learning is settled
per mission, not per round}: reusable state changes once, at the mission outcome.

\code{ReviseContract} is where evidence becomes an authorized change of target, and
its gate is deliberately two-tier. A \code{GoalContract} holds two kinds of clause.
Semantic clauses, exclusions, and recorded ambiguities move freely, because
clarifying what the user meant is ordinary Manager work. Precise clauses and the
objective itself require a confirmation covering exactly the clause identifiers that
change; without it the revision is refused. That split is the implemented form of the
distinction Section~\ref{sec:problem} draws analytically: intent is stable, wording
is negotiable, and the operational target moves only with recorded authority.

\subsection{Bounded Missions over Durable State}

The long-lived object in \systemname is the campaign, not a provider transcript. A
campaign has a persisted identity and objective; its work is divided into missions
that have explicit outcomes. The scheduler assigns one mission at a time, records
its result, and advances only at a clean mission boundary. Mission assignment is
transactional, preventing duplicate work under concurrency, and resumed execution
is tied to the persistent campaign identity.

Engineer and Reviewer calls use fresh provider sessions for each round. Cross-session
continuity comes from an ordinary shared \code{CHECKPOINT.md} containing durable
state, evidence references, open questions, and the next step. This follows the
broader move from monolithic context toward managed memory tiers and reusable
workflows~\citep{packer2023memgpt,wang2024workflowmemory,xu2025amem}. The Engineer
updates the checkpoint after execution. When independent review is invoked, the
Reviewer reads and corrects the same file and is its final editor for that round;
on an accepted self-review path, the Engineer's version remains the handoff. The
full history remains in artifacts and the event record.

This design makes restart and upgrade behavior part of the method. A running process
hands off only at a mission boundary, and the replacement resumes from the same
persistent campaign identity. After an interruption, replay or reassignment begins
from committed campaign state rather than reconstructing the task from a model
transcript.

\subsection{Review and Completion}

All role invocations pass through a common instrumentation layer that records usage
and associates each call with the mission trace. The Engineer and Reviewer share
the same artifact state, allowing review of the actual outputs and execution record
rather than only the Engineer's summary. A typed, append-only trace is the canonical
timeline; user interfaces are projections over that record.

The source of a mission completion verdict is explicit. When self-review is enabled
and no vertical or task policy requires independence, an Engineer may select
\code{review=skip}; its prompt restricts that choice to low-risk bounded work with
a passing verifier. The runtime accepts the explicit agent judgment without adding
a second heuristic or validator and records
\code{review\_source=engineer\_self\_review}. Otherwise a fresh Reviewer
returns the structured verdict, and stage-closing or vertical-required review cannot
be waived. Completion never comes from filenames, Token volume, or prose keywords.
The corresponding artifacts and execution records remain inspectable, while
credential redaction and result provenance preserve the selected verdict source.

\subsection{Verification-Gated Fixed-Model Runtime Self-Evolution}

Equation~\ref{eq:runtime-update} separates the model from the state around it.
Table~\ref{tab:runtime-ownership} identifies how each component can change and who
owns the committed update. The ownership model is intentionally not uniform. Memory
and skills use a work-versus-certification split: the Engineer or Scientist produces
a candidate, and the Reviewer commits the retained form. Tools and procedures are
system-configured. Stage checklists are Planner-owned, with Reviewer feedback rather
than a second commit gate. Routing policy is Manager-committed, while task
definitions are Planner-authored and scheduler-committed.

We use \emph{verification-guided} for the overall control policy that decides whether
to persist, stop, or pivot, and \emph{verification-gated} for the narrower admission
condition on reusable updates. We use \emph{runtime self-evolution} in a deliberately
narrow, fixed-model sense.
The underlying model parameters remain unchanged; what evolves is the persistent
state that later missions retrieve or obey. \emph{Verification-gated} is the
umbrella term used in this report: a
generated candidate is not reusable merely because a role produced it. Admission
requires the task-native evidence that exists for that surface and a commit by its
authorized owner. Depending on policy and risk, this can combine an official
executable verifier with independent Reviewer judgment, permitted Engineer
self-review, or a Manager, Planner, Scheduler, or system-configuration commit. It
does not mean that every low-risk update requires an independent Reviewer. We use
\emph{Reviewer-gated} only for the independent-Reviewer path and \emph{external
grader} only for a task-native evaluator outside the role loop.

A complete update cycle has four parts:
(1) an execution trajectory produces a candidate memory, skill, procedure,
verification rule, routing decision, or task definition; (2) the responsible role
checks the candidate against artifacts and task-native evidence; (3) the authorized
owner commits, revises, or rejects the update; and (4) a later mission retrieves the
retained state as part of its starting context or execution policy. Activity that
does not survive this commit-and-reuse path is not counted as self-evolution.

\begin{table}[H]
\centering
\small
\begin{tabularx}{\linewidth}{@{}p{2.25cm} p{2.75cm} p{2.55cm} X@{}}
\toprule
\rowcolor{papersand}
\textbf{State} & \textbf{Change source} & \textbf{Commit owner} & \textbf{Persistent form} \\
\midrule
$M$: memory & Engineer trajectory & Reviewer & Curated checkpoint and event-derived journal \\
$S$: skills & Engineer / Scientist & Reviewer & Versioned skill library \\
$A$: tools and procedures & System configuration & System configuration & Versioned tool and procedure registry \\
$V$: verification & Planner & Planner & Stage checklist; Reviewer supplies feedback \\
$R$: routing and roles & Runtime policy & Manager & Campaign routing policy \\
$Q$: tasks and evaluations & Planner & Scheduler & Backlog task specifications and scopes \\
\bottomrule
\end{tabularx}
\caption{Attributable ownership of fixed-model runtime evolution. A mission usually
updates only a subset of these components.}
\label{tab:runtime-ownership}
\end{table}

Two knowledge surfaces carry reusable experience. Versioned skills store procedures
that can be matched to later tasks. A project knowledge base stores source-linked
records and synthesized pages derived from reviewed outcomes. Neither surface is
treated as automatically correct: entries can be revised, archived, or retired when
later results contradict them.

This mechanism can improve a later mission without changing the model itself. A
verified retained failure can prevent a repeated dead end; an admitted procedure can shorten
environment inspection; a verifier can reject a previously accepted shortcut; and
a routing update can assign independent review to a higher-risk task. The mechanism
does not imply monotonic improvement. Some missions commit no reusable state,
retained state can become stale, and a harder task distribution can increase cost
even after useful state has accumulated.

\subsection{Reliability and Resource Governance}

Long-running systems fail through operational drift even when individual model
calls are strong. \systemname therefore combines round-level progress
classification with mission-boundary replanning. Work that repeatedly produces no
decision or artifact change is stopped or reformulated, while long-running external
jobs are tracked separately from model reasoning. These mechanisms bound execution;
scientific correctness remains with the task evaluator and the explicitly selected
Engineer-self-review or independent-Reviewer judgment path.

Resource accounting is centralized at model-call and external-job boundaries.
Budgets are checked before work begins and reconciled after completion, preventing
individual roles from expanding their own allocation. Concrete scheduling,
sandboxing, and deployment mechanisms are implementation choices rather than part
of the scientific claim.
